% Part: turing-machines
% Chapter: undecidability
% Section: representing-tms

\documentclass[../../../include/open-logic-section]{subfiles}

\begin{document}

\olfileid{tur}{und}{rep}
\olsection{টুরিং যন্ত্রের উপস্থাপন}

\begin{explain}
প্রথম-ক্রমের যুক্তিবিদ্যার !!a{sentence} দিয়ে টুরিং যন্ত্র ও তার আচরণ
উপস্থাপন করতে হলে আমাদের একটি উপযুক্ত ভাষা সংজ্ঞায়িত করতে হবে। ভাষাটির
দুটি অংশ: যন্ত্রের কনফিগারেশন বর্ণনার জন্য !!{predicate}s, আর নির্বাহের ধাপ
(“মুহূর্ত”) ও ফিতার অবস্থানগুলিকে সংখ্যা দেওয়ার রাশি।

আমরা দুই ধরনের !!{predicate}s প্রবর্তন করি, দু-ধরনেরই স্থানসংখ্যা~2: প্রতিটি
দশা~$q$-এর জন্য একটি !!a{predicate}~$\Obj Q_q$, এবং প্রতিটি ফিতা-প্রতীক
$\sigma$-এর জন্য একটি !!a{predicate}~$\Obj S_\sigma$। প্রথম ধরনের বিধেয়গুলি
$M$-এর দশা ও তার ফিতা-হেডের অবস্থান বর্ণনা করতে দেয়; দ্বিতীয় ধরনেরগুলি
ফিতার বিষয়বস্তু বর্ণনা করতে দেয়।

ফিতা-হেডের অবস্থান ও নির্বাহিত ধাপের সংখ্যা প্রকাশ করতে আমাদের সংখ্যা
প্রকাশের একটি উপায় চাই। এর জন্য ব্যবহার করা হয় একটি !!a{constant}~$\Obj 0$
এবং উত্তরসূরি অপেক্ষক নামে একটি $1$-স্থানীয় অপেক্ষক~$\prime$। রীতি অনুসারে এটি
তার আর্গুমেন্টের \emph{পরে} লেখা হয় (এবং বন্ধনী বাদ দেওয়া হয়)।

প্রতিটি সংখ্যা~$n$-এর একটি প্রামাণ্য পদ~$\num{n}$ আছে—$n$-এর
\emph{সংখ্যাপদ}—যা $\Lang L_M$-এ সংখ্যাটিকে উপস্থাপন করে। $\num{0}$ হলো
$\Obj 0$, $\num{1}$ হলো $\Obj 0'$, $\num{2}$ হলো $\Obj 0''$, ইত্যাদি। আরও
আনুষ্ঠানিকভাবে:
\begin{align*}
\num{0} & = \Obj 0 \\
\num{n+1} &= \num{n}'
\end{align*}
$\num{0}$ পদটি, অর্থাৎ $\Obj 0$, ফিতার বাঁদিকের প্রথম অবস্থানটির নাম দেয়;
একই সঙ্গে প্রথম নির্বাহ-ধাপের আগের সময়ের (আরম্ভিক কনফিগারেশনের) নামও দেয়।
$\num{1}$ পদটি, অর্থাৎ $\Obj 0'$, বাঁদিকের প্রথম ঘরের ডানের ঘরটির এবং প্রথম
নির্বাহ-ধাপের পরের সময়টির নাম দেয়; এভাবেই পরেরগুলি পাওয়া যায়।

ফিতার অবস্থানগুলির ক্রম (যেখানে এর মানে “বাঁ দিকে”) এবং নির্বাহের ধাপগুলির
ক্রম (যেখানে এর মানে “আগে”)—দুটিই প্রকাশ করতে আমরা !!a{predicate}~$<$-ও
প্রবর্তন করি।

ভাষাটি স্থির হয়ে গেলে আমরা $!T(M,w)$-এর “স্বতঃসিদ্ধগুলি”, অর্থাৎ একত্রে
$M$-কে নিবেশ~$w$-তে চালালে তার আচরণ বর্ণনা করে এমন !!{sentence}s তালিকাভুক্ত
করব। $\Obj 0$, $\prime$ ও~$<$-এর উপর শর্ত আরোপকারী !!{sentence}s, আরম্ভিক
কনফিগারেশন বর্ণনাকারী !!{sentence}s, এবং $M$ কোনো নির্দিষ্ট নির্দেশ পালন করার
পর তার কনফিগারেশন বর্ণনাকারী !!{sentence}s থাকবে।
\end{explain}

\begin{defn}
  \ollabel{defn:tm-descr}
একটি টুরিং যন্ত্র $M = \tuple{Q, \Sigma, q_0, \delta}$ দেওয়া থাকলে
ভাষা~$\Lang L_M$-এ থাকে:
\begin{enumerate}
\item প্রতিটি দশা~$q \in Q$-এর জন্য একটি দ্বি-স্থানীয় !!{predicate}
  $\Obj Q_q(x,y)$। স্বজ্ঞাতভাবে, $\Obj Q_q(\num{m},\num{n})$ প্রকাশ করে:
  “$n$ ধাপ পরে $M$ দশা~$q$-তে আছে এবং $m$-তম ঘরটি পড়ছে।”
\item প্রতিটি প্রতীক~$\sigma\in\Sigma$-এর জন্য একটি দ্বি-স্থানীয়
  !!{predicate} $\Obj S_\sigma(x,y)$। স্বজ্ঞাতভাবে,
  $\Obj S_\sigma(\num{m},\num{n})$ প্রকাশ করে: “$n$ ধাপ পরে $m$-তম ঘরে
  প্রতীক~$\sigma$ আছে।”
\item একটি !!{constant} $\Obj 0$
\item একটি এক-স্থানীয় !!{function} $\prime$
\item একটি দ্বি-স্থানীয় !!{predicate} $<$
\end{enumerate}
\end{defn}

$w=\sigma_{i_1}\dots\sigma_{i_k}$ নিবেশে টুরিং যন্ত্র~$M$-এর কাজ
বর্ণনাকারী !!{sentence}s নিচে দেওয়া হলো:
\begin{enumerate}
\item সংখ্যা ও~$<$ বর্ণনাকারী স্বতঃসিদ্ধ:
\begin{enumerate}
\item প্রতিটি সংখ্যা তার উত্তরসূরির চেয়ে ছোট—এ কথা বলা !!^a{sentence}:
\[
\lforall[x][x < x']
\]
\item $<$ যে সক্রম, তা নিশ্চিত করা !!^a{sentence}:
\[
\lforall[x][\lforall[y][\lforall[z][
      ((x < y \land y < z) \lif x < z)]]]
\]
\end{enumerate}
\item আরম্ভিক কনফিগারেশন বর্ণনাকারী স্বতঃসিদ্ধ:
\begin{enumerate}
\item $0$ ধাপ পরে—যন্ত্রটি শুরু হওয়ার আগে—$M$ আরম্ভিক দশা~$q_0$-তে এবং
  ঘর~$1$ পড়ছে:
\[
\Obj Q_{q_0}(\num{1}, \num{0})
\]
\item প্রথম $k+1$টি ঘরে যথাক্রমে $\TMendtape$,
  $\sigma_{i_1}$, \dots, $\sigma_{i_k}$ প্রতীক আছে:
\[
\Obj S_\TMendtape(\num{0}, \num{0}) \land
\Obj S_{\sigma_{i_1}}(\num{1}, \num{0}) \land
\dots \land
\Obj S_{\sigma_{i_k}}(\num{k}, \num{0})
\]
\item অন্য সব স্থানে ফিতা খালি:
\[
\lforall[x][(\num{k} < x \lif \Obj S_\TMblank(x, \num{0}))]
\]
\end{enumerate}
\item এক কনফিগারেশন থেকে পরেরটিতে অবস্থান্তর বর্ণনাকারী স্বতঃসিদ্ধ:

নিচের আলোচনায় $!A(x,y)$ বলতে $\sigma\in\Sigma$-এর প্রত্যেকটির জন্য
নিচের আকারের সব !!{formula}-র সংযোজন বোঝানো হবে:
% BN-SRC-178 TARGET-CORRECTION-BEGIN
\par\smallskip\noindent\textnormal{\small\textbf{উৎস-সংশোধন (BN-SRC-178):} হিমায়িত ইংরেজি পাঠ এখানে $x$ ও~$y$ মুক্ত থাকা রাশিগুলিকে “sentences” বলেছে। বাংলা পাঠে সঠিক ব্যাকরণগত শ্রেণি !!{formula} রাখা হয়েছে; $x,y$-এর জায়গায় বদ্ধ পদ বসালে সংশ্লিষ্ট বাক্যগুলি পাওয়া যায়।} % BN-SRC-178 TARGET-CORRECTION-END
\[
\lforall[z][
  (((z < x \lor x < z) \land \Obj S_\sigma(z, y))
  \lif \Obj S_\sigma(z, y'))]
\]
$!A(\num{m},\num{n})$ দিয়ে আমরা প্রকাশ করি: “ঘর~$m$ ছাড়া, $n+1$ ধাপের
পর ফিতাটি $n$ ধাপের পরের ফিতার মতোই আছে।”
\begin{enumerate}
\item \ollabel{rep-right} প্রতিটি নির্দেশ
$\delta(q_i,\sigma)=\tuple{q_j,\sigma',\TMright}$-এর জন্য !!{sentence}:
\begin{align*}
& \lforall[x][\lforall[y][(
   (\Obj Q_{q_i}(x, y) \land \Obj S_{\sigma}(x, y)) \lif {}]] \\
&\qquad   (\Obj Q_{q_j}(x', y') \land \Obj S_{\sigma'}(x, y') \land
!A(x, y)))
\end{align*}
এটি বলে: $y$ ধাপ পরে যন্ত্রটি যদি দশা~$q_i$-তে থেকে প্রতীক~$\sigma$-সহ
ঘর~$x$ পড়ে, তবে $y+1$ ধাপ পরে সেটি ঘর~$x+1$ পড়ছে, দশা~$q_j$-তে আছে,
ঘর~$x$-এ এখন $\sigma'$ আছে, এবং $x$ ছাড়া প্রতিটি ঘরে $y$ ধাপের পরের
একই প্রতীক আছে।

\item \ollabel{rep-left} প্রতিটি নির্দেশ
$\delta(q_i,\sigma)=\tuple{q_j,\sigma',\TMleft}$-এর জন্য !!{sentence}:
\begin{align*}
& \lforall[x][\lforall[y][
    ((\Obj Q_{q_i}(x', y) \land \Obj S_{\sigma}(x', y)) \lif {}]]\\
& \qquad   (\Obj Q_{q_j}(x, y') \land \Obj S_{\sigma'}(x', y') \land
!A(x', y))) \land {}\\
& \lforall[y][((\Obj Q_{q_i}(\num{0}, y) \land \Obj S_{\sigma}(\num{0},
    y)) \lif {}]\\
& \qquad (\Obj Q_{q_j}(\num{0}, y') \land \Obj S_{\sigma'}(\num{0},
  y') \land !A(\num{0}, y)))
\end{align*}
% BN-SRC-179 TARGET-CORRECTION-BEGIN
\par\smallskip\noindent\textnormal{\small\textbf{উৎস-সংশোধন (BN-SRC-179):} বাঁ-চলনের প্রথম শাখায় হিমায়িত সূত্রটি $!A(x,y)$ দিয়ে নতুন হেড-অবস্থান~$x$ বাদ দিয়েছিল, অথচ লেখা হয় আগের হেড-অবস্থান~$x'$-এ। বাংলা সূত্রে তাই $!A(x',y)$ রাখা হয়েছে; এতে ঠিক লেখা ঘরটিই অপরিবর্তনশীলতার শর্ত থেকে বাদ পড়ে।} % BN-SRC-179 TARGET-CORRECTION-END
এটি কীভাবে কাজ করে, একটু ভেবে দেখো। এবার আমরা “ঘর~$x$ পড়লে \dots” দিয়ে
শুরু করিনি; শুরু করেছি “ঘর~$x+1$ পড়লে \dots” দিয়ে। বাঁ দিকে চললে পরের
ধাপে যন্ত্রটি ঘর~$x$ পড়ে। কিন্তু যে ঘরটিতে লেখা হয় সেটি~$x+1$। বিয়োগ বা
পূর্বসূরি অপেক্ষক নেই বলেই আমরা এভাবে লিখি।

লক্ষ করো, $x+1$ আকারের সংখ্যাগুলি হলো $1$, $2$, \dots। তাই যন্ত্রটি যখন
ঘর~$0$ পড়ছে এবং বাঁ দিকে যাওয়ার কথা, তখনকার ঘটনাটি এতে ধরা পড়ে না
(যন্ত্রটি অবশ্য বাঁ দিকে যেতে পারে না—একই জায়গায় থাকে)। দ্বিতীয় সংযোজনটি
এই বিশেষ ঘটনাই সামলায়: $y$ ধাপ পরে যন্ত্রটি যদি দশা~$q_i$-তে ঘর~$0$ পড়ে
এবং ঘর~$0$-তে প্রতীক~$\sigma$ থাকে, তবে $y+1$ ধাপ পরে সেটি এখনও ঘর~$0$
পড়ছে, এখন দশা~$q_j$-তে আছে, ঘর~$0$-তে প্রতীক~$\sigma'$ আছে, আর ঘর~$0$
ছাড়া অন্য ঘরগুলিতে $y$ ধাপ পরে থাকা প্রতীকগুলিই আছে।
\item \ollabel{rep-stay} প্রতিটি নির্দেশ
$\delta(q_i,\sigma)=\tuple{q_j,\sigma',\TMstay}$-এর জন্য !!{sentence}:
\begin{align*}
& \lforall[x][\lforall[y][(
   (\Obj Q_{q_i}(x, y) \land \Obj S_{\sigma}(x, y)) \lif {}]] \\
&\qquad   (\Obj Q_{q_j}(x, y') \land \Obj S_{\sigma'}(x, y') \land
!A(x, y)))
\end{align*}
\end{enumerate}
\end{enumerate}
টুরিং যন্ত্র~$M$ ও নিবেশ~$w$-এর জন্য উপরের সব !!{sentence}s-এর সংযোজনকে
$!T(M,w)$ বলা হবে।

$M$ শেষ পর্যন্ত থামে—এটি প্রকাশ করতে আমাদের এমন একটি !!{sentence} খুঁজতে হবে
যা বলে: “কয়েকটি ধাপ পরে অবস্থান্তর অপেক্ষকটি অসংজ্ঞায়িত হবে।” যেসব
$\tuple{q,\sigma}$-এর জন্য $\delta(q,\sigma)$ অসংজ্ঞায়িত, সেগুলির সেট
$X$ হোক। তাহলে $!E(M,w)$ হবে নিচের !!{sentence}:
\[
\lexists[x][\lexists[y][(\bigvee_{\tuple{q, \sigma} \in
      X}(\Obj Q_q(x, y) \land \Obj S_\sigma(x, y)))]]
\]

নির্দিষ্ট থামার দশা~$h$-সহ টুরিং যন্ত্র ব্যবহার করলে আরও সহজ হয়: তখন
!!{sentence}~$!E(M,w)$
\[
\lexists[x][\lexists[y][\Obj Q_h(x, y)]]
\]
প্রকাশ করে যে যন্ত্রটি শেষ পর্যন্ত থামে।

\begin{prop}
\ollabel{prop:mlessk}
$m<k$ হলে $!T(M,w) \Entails \num{m}<\num{k}$।
\end{prop}

\begin{proof}
অনুশীলনী।
\end{proof}

\begin{prob}
\olref[tur][und][rep]{prop:mlessk} প্রমাণ করো।
(ইঙ্গিত: $k-m$-এর উপর আরোহ ব্যবহার করো।)
\end{prob}

\end{document}
